%======================= DOCUMENTATION ET LIVRAISON =========================
\chapter{Documentation et livraison}

\section{Une documentation produite en continu}

La documentation a accompagné le développement à chaque étape, plutôt que
d'être produite en fin de projet. Elle est regroupée dans le dossier
\texttt{docs/} du dépôt et couvre à la fois les aspects techniques,
fonctionnels et organisationnels. Le tableau~\ref{tab:docs} en présente les
principaux documents.

\begin{table}[H]
\centering
\caption{Documents de référence du dépôt}
\label{tab:docs}
\small
% Noms de fichiers composés avec \path : coupure autorisée aux tirets et
% points, ce qui évite tout débordement dans la colonne voisine.
\begin{tabularx}{\textwidth}{@{}p{6.2cm}L@{}}
\toprule
\textbf{Document} & \textbf{Contenu} \\
\midrule
\path{documentation-technique.md} & Architecture, structure des dossiers, installation, 24~endpoints avec exemples, schéma des entités, limites, pistes de déploiement. \\
\path{schema-base-de-donnees.md} & Schéma PostgreSQL cible (tables, contraintes, relations). \\
\path{scenario-demonstration-finale.md} & Script de démonstration en neuf étapes, texte oral et points techniques. \\
\path{checklist-demonstration-finale.md} & Checklists avant, pendant et après la démonstration ; erreurs courantes et solutions. \\
\path{livraison-finale-semaine-6.md} & Document de synthèse de livraison, prêt à être remis à l'encadrante. \\
\path{etat-avancement-semaine-6.md} & Journal détaillé, jour par jour, de la semaine~6. \\
\bottomrule
\end{tabularx}
\end{table}

\section{Documentation de l'API}

L'API est documentée de deux manières complémentaires : de façon
\textbf{interactive} via \textbf{Swagger/OpenAPI} (accessible à \texttt{/docs},
avec exécution directe des endpoints), et de façon \textbf{écrite} dans
\texttt{documentation-technique.md}, qui recense l'ensemble des endpoints avec
leur méthode, leur rôle et un exemple de réponse. Ce double niveau garantit
qu'un développeur reprenant le projet dispose à la fois d'une référence testable
et d'une référence rédigée.

\section{Préparation de la démonstration}

En vue de la présentation finale, deux documents dédiés ont été produits :

\begin{itemize}
  \item un \textbf{scénario de démonstration} détaillant, en neuf étapes, le
        parcours exact à montrer, avec pour chaque étape un texte oral court et
        les points techniques à mentionner ;
  \item une \textbf{checklist de démonstration} regroupant les vérifications
        avant, pendant et après la démo, un tableau des incidents courants
        (backend non lancé, port occupé, page blanche, cache navigateur,
        export CSV, erreur API) avec leur solution rapide, ainsi qu'un ordre
        recommandé de captures d'écran pour préparer des supports.
\end{itemize}

\section{Livraison finale}

La livraison finale prend la forme d'un document de synthèse
(\texttt{livraison-finale-semaine-6.md}) reprenant le périmètre livré, les
modules, l'état technique, les commandes de lancement, les endpoints
principaux, les limites honnêtes et les prochaines améliorations. L'ensemble du
code et de la documentation est versionné sur \textbf{GitHub}, chaque incrément
ayant été intégré par \emph{pull request}, ce qui rend la livraison
\textbf{traçable} et \textbf{reproductible}.
